وتشير التوقعات إلى استمرار التراجع الطفيف في معدل التضخم العالمي من 4.1\% في عام 2025 إلى نحو 3.8\% في عام 2026، مدفوعا بتراجع الطلب وتراجع أسعار العديد من السلع الأساسية عالميا. إلا أن مسار التضخم يظل متفاوتا بين الدول؛ إذ شهدت الاقتصادات المتقدمة بما فيها الولايات المتحدة والاتحاد الأوروبي تراجعا ملحوظا في معدلات التضخم نتيجة تشديد السياسات النقدية، بينما ظلت ديناميكيات التضخم في الأسواق الناشئة متباينة، مما يعكس العوامل والمحركات الخاصة بكل دولة. ورغم ذلك، حافظت البنوك المركزية الكبرى إلى حد كبير على نهج حذر في تعديل سياساتها النقدية.
